\chapter{FOUNDATIONS OF MATHEMATICS; LOGIC; SET THEORY.}

The study of what is usually called the ``foundations of mathematics'', which has been carried out ceaselessly since the beginning of the 19th century, was impossible to bring to fruition except with the help of a parallel effort to systematise Logic, at least those parts that govern the links between mathematical statements. So it is not possible to separate the history of set theory and the formalisation of mathematics from that of ``Mathematical Logic''. But traditional logic, like that of the modern philosophers, covers in principle, an area of applications far greater than Mathematics. Therefore the reader must not expect to find in what follows a history of Logic, even in a very summarised form; we have limited ourselves as far as possible to retracing the evolution of Logic only in so far as it impinged on that of Mathematics. It is because of this that we will say nothing about the non-classical logics (many-valued logics, modal logics); all the more so will we be unable to tackle the history of those controversies which, from the Sophists to the Vienna School, have never stopped dividing philosophers both as to the possibility and the manner of applying Logic to objects in the real world or to concepts of human thought.

That there was a well-developed prehellinic mathematics is not today in any doubt. Not only are the notions (already very abstract) of whole number and of the measurement of quantities commonly used in the most ancient documents which have reached us from Egypt or Chaldea, but Babylonian algebra, because of the elegance and sureness of its methods, should not be thought of as a simple collection of problems solved by empirical fumbling. And, if nothing is found in the texts which resembles a ``proof'' in the formal meaning of the word, it is reasonable to believe that the discovery of such methods of solution, whose generality appears through particular numerical applications, was not possible without a minimum of logical links (perhaps not entirely conscious, but rather like those on which a modern algebraist depends when he undertakes a calculation, before ``setting down formally'' all the details) ({[}232{]}, pp.~203 ff.).

The essential originality of the Greeks consisted precisely of a conscious effort to order mathematical proofs in a sequence such that passing from one link to the next leaves no room for doubt and constrains universal assent.

That Greek mathematicians made use, in the course of their research, just like modern mathematicians, of ``heuristic'', rather than convincing, arguments, is what was proved for example (if there were any need), by the ``treatise of method'' of Archimedes {[}153 c{]}; note also in this, allusions to results ``found, but not proven'' by earlier mathematicians. \(^{1}\) But from the first detailed texts that are known to us (and which date from the middle of the Vth century), the ideal ``canon'' of a mathematical text is properly settled. It will find its highest expression in the great classics, Euclid, Archimedes and Apollonius; the notion of proof, in these authors, differs in no way from ours.

We have no texts allowing us to follow the first steps of this ``deductive method'', which seems to us already near perfection at the exact moment when we become aware of its existence. One can only think that it fits fairly naturally into the perpetual search for ``explanations'' of the world which characterises Greek thought and which is so discernible already amongst the Ionian philosophers of the VIIth century; moreover tradition is unanimous in ascribing the development and refinement of the method to the Pythagorean School, in a period which fits between the end of the VIth century and the middle of the Vth century.

It is on this ``deductive'' mathematics, fully conscious of its goals and methods, that the philosophical and mathematical thought of subsequent times will be concentrated. We will see on the one hand the establishment little by little of ``formal'' Logic modelled on mathematics, to conclude with the creation of formalised languages; on the other hand, mainly starting at the beginning of the XIXth century, the basic concepts of mathematics will be queried more and more and a great effort will be made to clarify their nature, especially after the birth of the Theory of sets.

